\title{Parameter-Efficient Orthogonal Finetuning via Factorization}

\begin{abstract}
The widespread adoption of large foundation models has made training them from the ground up highly cost-prohibitive. As such, effectively adapting these large-scale models to specific downstream tasks is becoming increasingly critical. This work focuses on Orthogonal Finetuning (OFT), a well-founded approach for downstream task adaptation. While OFT demonstrates robust generalization properties, it still requires training a substantial number of parameters, attributable to the significant dimensionality inherent in orthogonal matrices. To enhance parameter-efficiency, we first re-examine OFT through the lens of information transmission and distill several essential criteria for improved parameter usage. Drawing inspiration from the Cooley-Tukey fast Fourier transform's efficient handling of information flow, we devise a compact orthogonal parameterization leveraging butterfly structures. Integrating this parameterization into OFT, we introduce Orthogonal Butterfly (BOFT), a novel, parameter-efficient finetuning technique. BOFT generalizes the OFT framework, effectively subsuming it as a specific instance. We further provide a comprehensive empirical evaluation, demonstrating BOFT's effectiveness for adapting large vision transformers, expansive language models, and text-to-image diffusion models across a spectrum of vision and language downstream applications.
\end{abstract}

\section{Introduction}

Modern models like ChatGPT~\cite{brown2020language,chen2021evaluating} and Stable Diffusion~\cite{rombach2022high} showcase the exceptional generalization performance characteristic of large-scale foundation models. This surge in capability is mirrored by a steep increase in model size

(\emph{e.g.}, GPT-3 features approximately 175B parameters), making it increasingly difficult for the research community to train such foundational architectures from scratch. To advance the field, it is crucial to develop techniques that facilitate the adaptation of pre-trained foundation models to new tasks without the computational demands of complete retraining. Effective task adaptation strategies fall into three principal categories: \emph{model finetuning}~\cite{hulora2022,zaken2022bitfit,guo2020parameter,raffel2020exploring,qiu2023controlling,zhang2022adaptive,chavan2023one}, which involves adjusting only a subset of the existing model's parameters while leaving the architecture unaltered; \emph{adapter tuning}~\cite{rebuffi2017learning,houlsby2019parameter,pfeiffer2020adapterfusion,lin2020exploring,he2021towards}, which integrates additional trainable components into the model; and \emph{prompt tuning}~\cite{li2021prefix,lester2021power}, where extra prefix tokens are learned and attached to the input data. Among these, model finetuning stands out due to its simplicity and efficacy, as well as its advantage of not incurring additional inference delays.

Model finetuning fundamentally seeks to maintain similarity between the pretrained and finetuned models according to specific metrics, thereby retaining the knowledge acquired during pretraining. Typically, this is achieved by employing a low learning rate in most conventional finetuning strategies, a pragmatic technique that limits the extent of divergence between the original and finetuned model parameters. Recognizing the inherent structure within weight matrices, a more theoretically motivated strategy aims to conserve the relational properties of these matrices—specifically, by preserving the pairwise angular relationships among neurons. Building on this premise, the Orthogonal Finetuning~(OFT)~\cite{qiu2023controlling} framework was proposed, in which all neurons within a single layer undergo a transformation via a shared orthogonal matrix, guaranteeing that pairwise angles between neurons remain constant throughout finetuning. Although OFT has shown substantial improvements in generalizability and convergence for tasks such as finetuning text-to-image diffusion models~\cite{qiu2023controlling}, it suffers from the challenge that orthogonal matrices entail a large number of parameters due to their high dimensionality. To alleviate this, OFT employs a block-diagonal form, thereby curtailing the parameter count. Nonetheless, this reduction in parameters brings its own limitations—the orthogonal matrices have a predetermined sparsity configuration, and orthogonal transformations act independently across distinct blocks. While this block-diagonal approach yields parameter savings, it introduces potential inductive biases (\emph{e.g.}, diminishing the expressive capacity of the transformations and making them unsuitable for approximating traditional linear operations). Moreover, the optimal strategy for selecting a block-wise sparsity pattern is still an open problem.

A central challenge in this context is the creation of dense orthogonal matrices without incurring excessive parameter counts. Traditionally, forming a dense orthogonal matrix in $d$ dimensions necessitates $\mathcal{O}(d^2)$ parameters, which poses practical limitations. Our strategy circumvents this by assembling dense orthogonal matrices through the composition of several factorized sparse matrices. The rationale is to substitute increased computational overhead for a reduced number of tunable parameters: because the orthogonal matrix is expressed as a product of sparse factors, these multiplications are executed multiple times during training. Conceptually, we recast this factorization as an information transmission scenario, where generating a dense orthogonal matrix via products of sparse matrices mirrors information flow across a grid-shaped graph. This lens not only clarifies the factorization process but also supports the design of multiple parameter-efficient schemes, where the expressive power needed to generate dense matrices is retained while the count of learnable parameters remains low.

To advance parameter efficiency within our information transmission paradigm, we draw upon the butterfly architecture underpinning the Cooley-Tukey fast Fourier transform algorithm~\cite{cooley1965algorithm}, a topology known for its effective communication between graph nodes~\cite{klasing1994broadcasting}. The butterfly network embedded within the Cooley-Tukey method leads naturally to a form of sparse matrix decomposition—termed butterfly factorization. With butterfly factorization, a dense $d$-dimensional matrix emerges as a product of $\mathcal{O}(\log d)$ sparse matrices, each containing only $\mathcal{O}(d)$ non-zero entries. By maintaining orthogonality at each stage, we obtain an efficient parameterization for orthogonal matrices with merely $\mathcal{O}(d\log d)$ parameters, representing a substantial saving relative to the naive approach. Leveraging this compact orthogonal parameterization, we introduce Orthogonal Butterfly (BOFT), a new finetuning methodology that generalizes OFT, thus establishing a broad orthogonal finetuning framework. Notably, both the block-diagonal and butterfly constructions exploit matrix sparsity as a means to minimize the number of trainable parameters. In comparison, the low-rank structures featured in LoRA~\cite{hulora2022} offer a parallel form of efficiency; in fact, both low-rank and sparse matrices belong to the primary classes of structured matrices that facilitate parameter-efficient optimization~\cite{candes2011robust}.

In contrast to OFT’s reliance on a block-diagonal configuration to balance model expressivity against regularization, BOFT employs the butterfly architecture to more gradually interpolate between matrices from the full orthogonal group (emphasizing expressivity) and the identity matrix (encouraging regularity). This facilitates the isolation of a more compact hypothesis space within the orthogonal group suitable for downstream applications. Given the centrality of butterfly structures in accelerating classical linear operations such as the discrete Fourier and discrete cosine transforms, we posit that embedding such structure into parameter-efficient finetuning injects a helpful inductive bias—one that not only enhances generalization but also guards against overfitting. Our reasoning is rooted in the butterfly structure’s capacity to exactly represent a broad array of fundamental linear transformations—a feat unattainable by the block-diagonal parameterization used in OFT. Our main contributions are summarized as follows:

\begin{itemize}[leftmargin=*,nosep]
    \item We address the challenge of parameter efficiency in orthogonal finetuning via a novel framework grounded in information transmission. This perspective recasts the pursuit of dense, parameter-efficient orthogonal matrices as the analysis of information flow on grid-structured graphs.
    \item Drawing inspiration from the butterfly factorization intrinsic to the Cooley-Tukey algorithm, we introduce Orthogonal Butterfly—a new, parameter-efficient orthogonal finetuning technique situated within the information transmission framework.
    \item We offer theoretical analysis to elucidate how BOFT achieves significant reductions in the quantity of trainable parameters without sacrificing either expressivity or the stability of optimization. Moreover, the underlying matrix factorization enables an effective form of weight interpolation (see Figure~\ref{fig:boft_interpolation}).
    \item For the first time, we systematically apply orthogonal finetuning to a broad range of tasks that extend beyond controllable text-to-image generation~\cite{qiu2023controlling}, underscoring its promise as a universal approach to model adaptation. Specifically, we demonstrate BOFT’s utility across diverse adaptation scenarios, spanning computer vision and natural language processing. Our results show that BOFT consistently surpasses leading methods, thereby substantiating its advantages in parameter efficiency and generalization.
\end{itemize}

\textbf{Parameter-efficient finetuning (PEFT)}. The rapid expansion and increased capabilities of foundation models (\emph{e.g.}, \cite{brown2020language,radford2021learning,rombach2022high,kirillov2023segment}) have elevated interest in strategies that efficiently finetune these models for downstream applications while minimizing parameter usage~\cite{treviso2023efficient,ding2023parameter,lialin2023scaling}. Among the many PEFT strategies introduced~\cite{houlsby2019parameter,aghajanyan2020intrinsic,hulora2022,edalati2022krona,wang2022adamix,gheini2021cross,zaken2022bitfit,guo2020parameter,sung2021training,ansell2022composable,lester2021power,li2021prefix,vu2022spot,he2021towards,mao2021unipelt,karimi2021compacter,liu2022few,sung2022lst,chen2022parameter,jia2022visual,chen2022adaptformer,zhang2022neural,jie2023fact,lian2022scaling,luo2023towards}, approaches based on reparameterization~\cite{aghajanyan2020intrinsic,hulora2022,edalati2022krona} are especially pertinent to our research. LoRA~\cite{hulora2022}, for instance, enhances a pretrained weight matrix by incorporating the product of two low-rank matrices, thereby obtaining strong outcomes on various natural language tasks. Unlike LoRA’s use of a fixed matrix rank across all layers, some recent methods~\cite{zhang2022adaptive,valipour2022dylora,zhang2023increlora} dynamically adjust the rank per layer, achieving a more adaptive allocation of the parameter budget. The straightforward nature of low-rank weight reparameterization has contributed to its wide adoption~\cite{dettmers2023qlora,zi2023delta,chavan2023one}. Inspired by the application of hyperspherical energy to generalization analysis~\cite{liu2018learning,liu2021orthogonal}, \cite{qiu2023controlling} introduces orthogonal finetuning (OFT), an alternative method specifically designed for finetuning text-to-image diffusion models. In OFT, an orthogonal matrix is learned to transform the neurons within a given layer, producing stronger generalization capabilities and more stable training behaviors when compared to LoRA. Nevertheless, OFT typically requires a greater number of trainable parameters than LoRA, making improvements in its parameter efficiency desirable. Furthermore, the effectiveness of OFT in contexts beyond text-to-image diffusion model adaptation~\cite{qiu2023controlling} has not been thoroughly explored. The introduction of BOFT, leveraging butterfly factorization, enhances the parameter efficiency of OFT. This advancement enables us to empirically validate the effectiveness of orthogonal finetuning in a broader range of adaptation tasks.

\textbf{Butterfly structures}. The radix-2 Cooley-Tukey algorithm decomposes the $N$-point discrete Fourier transform into two separate $\frac{N}{2}$-point transforms in a recursive manner, a process that inherently generates a butterfly structure. This structure is characterized by its formulation as a product of several sparse matrices, collectively known as a butterfly matrix. Such matrices have a history of usage in parameterizing orthogonal matrices to simplify computations like Gaussian elimination without the need for pivoting and to boost computational efficiency~\cite{parker1995random}. They have also proven beneficial for stabilizing the training dynamics of recurrent neural networks~\cite{jing2017tunable} and for kernel approximation~\cite{munkhoeva2018quadrature}. Methods for learning rapid linear transforms via butterfly parameterization have been described in \cite{dao2019learning,dao2019kaleidoscope}. In neural network training, butterfly matrices facilitate parameter-efficient architectures~\cite{chen2022pixelated,dao2022monarch}. Additional applications include accelerating matrix-vector products~\cite{michielssen1996multilevel,o2010algorithm}, compressing matrices with data-sparse approximations~\cite{li2015butterfly}, and enabling efficient network transmission~\cite{klasing1994broadcasting,li2003linear,rajasingh2016transmission}. Our work diverges from these prior efforts by specifically leveraging butterfly structures to advance the parameter efficiency of OFT.

\section{An Information Transmission View on Orthogonal Finetuning}

We begin by reviewing some foundational concepts of OFT. In order to finetune a pretrained weight matrix, OFT expresses the updated weight as the product of a trainable orthogonal matrix and the original, fixed pretrained weight matrix. Unlike LoRA, which introduces an additive low-rank matrix for weight adaptation, OFT incorporates a multiplicative orthogonal transformation. LoRA’s parameter efficiency is achieved through its low-rank formulation, whereas the standard OFT attains efficiency by employing a block-diagonal orthogonal structure~\cite{qiu2023controlling}; parameter usage decreases as the block size along the diagonal shrinks. Figure~\ref{fig:comp_lora} illustrates a straightforward comparison. The rationale for utilizing orthogonal transformations in finetuning lies in their ability to maintain pairwise angular relationships among neurons~\cite{liu2018learning,liu2021orthogonal,qiu2023controlling}, which facilitates the retention of semantic knowledge encoded during pretraining. Specifically, OFT learns an orthogonal matrix $\bm{R}\in\mathbb{R}^{d\times d}$ for a pretrained linear transformation $\bm{W}^0\in\mathbb{R}^{d\times n}$, modifying the original forward computation from $\bm{z}=(\bm{W}^0)^\top\bm{x}$ to $\bm{z}=(\bm{R}\bm{W}^0)^\top\bm{x}$, where $\bm{x}\in\mathbb{R}^{d}$ denotes the input and $\bm{z}\in\mathbb{R}^{n}$ is the output. For zero initialization, $\bm{R}$ is set as the identity matrix at the start. To guarantee that $\bm{R}$ remains orthogonal during finetuning, the Cayley parameterization is applied as in \cite{qiu2023controlling,liu2021orthogonal}: $\bm{R}=(\bm{I}+\bm{Q})(\bm{I}-\bm{Q})^{-1}$, with $\bm{Q}$ being a skew-symmetric matrix, i.e., $\bm{Q}=-\bm{Q}^\top$. For improved parameter efficiency, the orthogonal matrix $\bm{R}$ adopts a block-diagonal form, parameterized as $\text{diag}(\bm{R}_1,\bm{R}_2,\cdots,\bm{R}_r)$, where each block $\bm{R}_i\in\mathbb{R}^{b\times b}$ is itself orthogonal and $br=d$. This block-diagonal design streamlines parameter count, but it imposes a constraint: it effectively partitions each neuron's dimensions (columns of $\bm{W}^0$) into $r$ groups, each being independently transformed via separate orthogonal matrices. While this structure proves effective empirically, there is no theoretical justification for dividing neuron dimensions solely based on index grouping, highlighting the appeal of employing fully dense orthogonal matrices. Thus, we are prompted to ask: \emph{Is it feasible to devise a dense orthogonal matrix that also remains parameter-efficient?}

To investigate this issue, we suggest expressing a dense orthogonal matrix as the product of several sparse orthogonal matrices. In order to systematically and intuitively analyze the sparsity structure in orthogonal matrix factorizations, we reinterpret the task of generating a dense orthogonal matrix within OFT through the lens of information transmission. More precisely, forming a dense matrix $\bm{R}\in\mathbb{R}^{d\times d}$ by multiplying $m$ square matrices as $\bm{R}=\bm{B}_m\bm{B}_{m-1}\cdots\bm{B}_1$ can be conceptualized as the movement of information across a grid consisting of $d\times (m+1)$ nodes, as depicted in Figure~\ref{fig:info_trans}. This perspective is motivated by recognizing that a dense $d$-dimensional square matrix naturally represents complete connectivity from one group of $d$ nodes to another group of $d$ nodes. For the given matrix $\bm{R}$, if the entry $R_{ij}$ equals zero, it signifies the inability of information to propagate from the $j$-th source node to the $i$-th target node; conversely, a nonzero $R_{ij}$ allows information transfer. Therefore, decomposing $\bm{R}$ into the product $\bm{B}_m\bm{B}_{m-1}\cdots\bm{B}_1$ corresponds to a sequential information passing process guided by the structure of the graphs defined by each $\bm{B}_i$. The flow of information commences with $\bm{B}_1$ and concludes with $\bm{B}_n$. As an illustrative case presented in Figure~\ref{fig:info_trans}, consider the factorization $\bm{R}=\bm{B}_5\bm{B}_4\bm{B}_3\bm{B}_2\bm{B}_1$, for which both the sparsity patterns and their associated graphs are visualized. The depicted graph arises from unfolding the matrix multiplication process: each matrix $\bm{B}_i$ provides a connectivity map linking nodes from level $i$ to those in level $i+1$. Specifically, the entry $(j_1, j_2)$ in $\bm{B}_i$ indicates the existence (or absence, if zero) of a directed connection from the $j_2$-th node at level $i$ to the $j_1$-th node at level $i+1$. For the full product $\bm{B}_5\bm{B}_4\bm{B}_3\bm{B}_2\bm{B}_1$ to yield a dense matrix, it is necessary that every initial node can transmit information to all terminal nodes at the sixth level. If we examine the case $\bm{R}=\bm{B}_4\bm{B}_3\bm{B}_2\bm{B}_1$—corresponding to connections from the first to the fifth level—we observe, for example, that information originating at node 1 cannot reach node 3 at the endpoint; thus, the sparsity pattern of $\bm{B}_4\bm{B}_3\bm{B}_2\bm{B}_1$ includes a zero in position $(3,1)$. This correspondence between graph structure and sparsity extends to shorter products, such as $\bm{R}=\bm{B}_3\bm{B}_2\bm{B}_1$ and $\bm{R}=\bm{B}_2\bm{B}_1$. In general, for a matrix $\bm{R}\in\mathbb{R}^{d\times d}$ to be dense, the sequence of $m$ matrices $\bm{B}_m, \dots, \bm{B}_1$ should together encode a set of directed edges in a $d\times(m+1)$ grid where each edge links only adjacent levels (i.e., adjacent columns); this guarantees that any node from the initial level can transfer information to every node in the terminal $(m+1)$-th level.

The information transmission perspective offers valuable insights into the sparsity characteristics found in the factorization matrices within OFT. As depicted in Figure~\ref{fig:blk_diag}, the matrix $\bm{R}$ in the standard OFT exhibits a block-diagonal configuration. Although this arrangement serves to lower the count of trainable parameters, it falls short of producing a fully dense $\bm{R}$. Our objective is to assemble a dense orthogonal matrix from $m$ sparse orthogonal matrices, optimizing for the least possible number of trainable parameters. Within the information transmission framework, the essential requirements for achieving this objective are: \emph{(i)} \textbf{dense connectivity}, ensuring that every node at the input layer is connected through some sequence of paths to all nodes at the output layer, and \emph{(ii)} \textbf{minimal free edges}, such that the total edge count is minimized while respecting the orthogonality property. It is important to recognize that orthogonality enforces nuanced restrictions on the connections (edges) between consecutive layers. In particular, maintaining orthogonality necessitates each matrix $\bm{B}_i$ to be full-rank, requiring $d$ connections forming a bijective correspondence between nodes at level $i$ and those at level $(i+1)$. This means there must be at least $d$ connections between successive layers (for instance, 4 edges in the scenario illustrated in Figure~\ref{fig:info_trans}). These $d$ required edges ensure orthogonality and do not count towards the tally of edges relevant for trainability, since these elements remain fixed and non-trainable (e.g., for a $d\times d$ orthogonal matrix with $d$ nonzero components, they can only take values $\pm 1$). Given that orthogonal matrices use fewer parameters than generic matrices, the orthogonality condition introduces further interdependencies among the edges. As a specific case, in a $2\times 2$ orthogonal matrix, a zero at entry $(1,1)$ implies another zero at $(2,2)$—that is, the absence of one edge necessitates the absence of another. Consequently, for each permissible configuration of edges, the orthogonality constraint may sometimes necessitate the addition or removal of connections. When examining the structure of non-zero entries in the resulting orthogonal matrix, the information transmission paradigm becomes particularly advantageous in OFT, as the focus lies exclusively on the positions of trainable nonzero elements in $\bm{R}$, and their actual values are inconsequential.

A straightforward dense linkage between two hierarchical layers requires $\mathcal{O}(d^2)$ connections (essentially equivalent to a dense orthogonal matrix), which results in $d^2-d$ connections; for instance, with $d=4$, this amounts to 12 connections. Figure~\ref{fig:info_trans} demonstrates an example utilizing matrix factorization, achieving the required connectivity with just $10$ edges—fewer than that needed for a fully dense orthogonal matrix. This framework allows us to investigate the parameter-efficiency of OFT through a graphical lens and facilitates the design of feasible matrix factorizations. Drawing from a compelling network architecture known as the butterfly graph~\cite{cooley1965algorithm}, as seen in the Cooley-Tukey algorithm, we observe that it provides an efficient mechanism to connect $d$ input nodes to $d$ output nodes with only $\mathcal{O}(d\log d)$ edges. For example, whereas the configuration in Figure~\ref{fig:info_trans} requires $10$ edges for full connectivity, the butterfly network accomplishes this with merely $8$. In the next section, we elaborate on how adopting the butterfly structure enhances parameter efficiency.

\section{Orthogonal Parameterization by Butterfly Factorization}

The butterfly network originates from the Cooley-Tukey algorithm for efficient computation of the fast Fourier transform. In the Fourier transform context, localized adjustments in the frequency domain correspond to widespread changes in the spatial domain, aligning well with our scenario of information dissemination—every node at the input layer is capable of transmitting information to every node at the output layer. The butterfly topology has also been widely adopted in computer networking~\cite{solihin2015fundamentals,li2011linear} for efficient data transfer. Let $k\geq 2$ denote a power of two; we formally introduce the \emph{butterfly factor} $\bm{B}^F(k)$ as follows:

\begin{equation}\label{eq:but_factor}
\begin{aligned}
    \bm{B}^F(k)=\begin{bmatrix}
    \text{diag}(\bm{d}_1) & \text{diag}(\bm{d}_2) \\
    \text{diag}(\bm{d}_3) & \text{diag}(\bm{d}_4)
    \end{bmatrix}\in\mathbb{R}^{k\times k},
\end{aligned}
\end{equation}

where each vector $\bm{d}_i \in \mathbb{R}^{\frac{k}{2}}$ for $i=1,2,3,4$. For $d=2^N$, the \emph{butterfly matrix} $\bm{B}(d)\in\mathbb{R}^{d\times d}$ is constructed recursively as

\begin{equation}
\begin{aligned}
    \!\!\bm{B}(d)=\tilde{\bm{B}}(d,d)\!\cdot\!\begin{bmatrix}
    \bm{B}_1(\frac{d}{2})\!\!\!\!\! & \bm{0} \\
    \bm{0}\!\!\!\!\! &\bm{B}_2(\frac{d}{2})
    \end{bmatrix}= \tilde{\bm{B}}(d,d)\tilde{\bm{B}}(d,\frac{d}{2})\cdots\tilde{\bm{B}}(d,2),
\end{aligned}
\end{equation}

where $\bm{B}_1(\frac{d}{2})$ and $\bm{B}_2(\frac{d}{2})$ denote two butterfly matrices of dimension $\frac{d}{2}$. We introduce the notion of a \emph{butterfly component} as $\thickmuskip=2mu \medmuskip=2mu \tilde{\bm{B}}(d,k)=\text{diag}(\bm{B}^F_1(k),\cdots,\bm{B}^F_{d/k}(k))$, which represents a block-diagonal matrix of size $d\times d$ whose blocks are each of size $k$. Each $\bm{B}^F_i(k),\forall i$ corresponds to a butterfly factor, as given in Equation~\ref{eq:but_factor}. Utilizing the butterfly matrix structure, we seek to parameterize an orthogonal matrix. For this construction to maintain orthogonality, all product factors $\tilde{\bm{B}}(d,k),\forall k$ composing $\bm{B}(d)$ must be orthogonal matrices. We begin by examining $\tilde{\bm{B}}(d,2)$, the block-diagonal matrix with $2\times2$ blocks. Each $2\times2$ block can be parameterized via the Cayley transform (equivalent to a $2$-dimensional rotation), ensuring orthogonality, as previously established in \cite{liu2021orthogonal,qiu2023controlling}. The sparsity structure for every butterfly component corresponds to a permutation of the nonzero arrangement in $\tilde{\bm{B}}(d,2)$, allowing all components to be systematically parameterized as orthogonal. This approach leads to a computationally efficient scheme for representing orthogonal matrices using collections of $2\times2$ orthogonal blocks. Drawing on the methodology in \cite{chen2022pixelated}, we extend the butterfly matrix concept to define a block butterfly component $\tilde{\bm{B}}^b(d,k)$, where every entry in $\bm{d}_i$ is replaced by a $b\times b$ matrix. For the block butterfly component $\tilde{\bm{B}}^b(d,2)$, orthogonality is enforced by parameterizing each $2b\times 2b$ block as an orthogonal matrix. For components with $k>2$, i.e., $\tilde{\bm{B}}^b(d,k)$, the pattern of nonzero blocks is simply a permuted block-wise arrangement of the sparsity in $\tilde{\bm{B}}^b(d,2)$, which means these too can be parameterized as orthogonal matrices in an analogous manner. Collecting these elements, the BOFT forward pass can be expressed as

\begin{equation}
    \begin{aligned}\nonumber
    \bm{z}=\big(\bm{R}(m,b)\cdot\bm{W}^0\big)^\top\bm{x},~~~~\text{s.t.}~\bigg{\{}\bm{R}(m,b)=\prod_{i=1}^m \tilde{\bm{B}}^b_{(d,i)}~~\&~~\underbrace{\big(\tilde{\bm{B}}^b_{(d,j)}\big)^\top\tilde{\bm{B}}^b_{(d,j)}=\tilde{\bm{B}}^b_{(d,j)}\big(\tilde{\bm{B}}^b_{(d,j)}\big)^\top\!=\bm{I}_d}_{\forall j\in [1, m]}\bigg{\}},
    \end{aligned}
\end{equation}

For notational convenience, we abbreviate $\tilde{\bm{B}}^b(d,2^{m - i+1})$ as $\tilde{\bm{B}}^b_{(d,i)}$, with $\bm{I}_d$ denoting the $d$-dimensional identity matrix. The orthogonal matrix $\bm{R}(m,b)\in\mathbb{R}^{d\times d}$ is constructed via a product of $m$ orthogonal butterfly factors. For ease of reference, we represent BOFT using $\bm{R}(m,\frac{b}{2})$ as BOFT($m$,$b$), provided $b\geq 2$. When $m=1$, BOFT($1$,$b$) specializes to the block-diagonal OFT~\cite{qiu2023controlling}, featuring blocks of dimension $b$. Setting $b=d$, BOFT($1$,$d$) simplifies to the conventional OFT~\cite{qiu2023controlling}, wherein a fully unrestricted orthogonal matrix is employed. With the choice BOFT($\log \frac{2d}{b}$,$b$), the construction utilizes the block butterfly matrix $\bm{B}^{\frac{b}{2}}(d)$ for $\bm{R}$, thereby producing a dense orthogonal structure. Broadly, BOFT($m$,$b$) introduces $\frac{1}{2}(b-1)dm$ adjustable parameters that can be finetuned for a linear transformation of dimension $d\times n$. If a standard butterfly parameterization is utilized, i.e., $m=\log d$, $b=2$, BOFT requires only $\mathcal{O}(d\log d)$ parameters. By contrast, the full dense OFT necessitates $\mathcal{O}(d^2)$ parameters, while its block-diagonal variant, using $r$ blocks, requires $\mathcal{O}(bd)$. Notably, whereas the traditional OFT can only achieve a dense matrix by choosing $b=d$, BOFT achieves a dense orthogonal parameterization with any block size $b$.

\textbf{Identity initialization for BOFT}. Typically, finetuning procedures commence from the parameters of a pretrained network to ensure that the optimization process starts in a neighborhood of the original model. As an illustration, LoRA initializes its low-rank updates with zeros. In the context of BOFT, every butterfly factor is initialized as an identity matrix (specifically, the skew-symmetric part for the Cayley transform is set to zero).

\textbf{Multiplicative Dropout for BOFT}. LoRA~\cite{hulora2022} adopts Dropout regularization for its low-rank updates to alleviate overfitting. While standard Dropout~\cite{srivastava2014dropout} integrates naturally with LoRA, it is incompatible with BOFT’s multiplicative update mechanism. To reconcile this, we introduce a form of multiplicative Dropout specifically designed for BOFT. Given that $\bm{R}(m,b)$ comprises $m$ butterfly layers, each of which can be decomposed into $2b\times 2b$ block-diagonal orthogonal components, the multiplicative Dropout operates by stochastically selecting $p_1$ proportion of butterfly modules and $p_2$ proportion of block-diagonal segments per module, subsequently replacing those selected blocks with identity matrices.

\section{Intriguing Insights and Discussions}

\textbf{Expressivity of BOFT}. While the butterfly structure combined with permutations is known to exactly realize various well-established fast linear transforms~\cite{dao2019learning,dao2019kaleidoscope} (\emph{e.g.}, fast Fourier transform, Hadamard transform), the extent to which our orthogonal butterfly matrix can approximate a general orthogonal matrix has not been fully characterized. To address this, we perform an empirical study targeting the approximation of a randomly generated dense orthogonal matrix~\cite{anderson1987generation} of dimension $512\times 512$. Figure~\ref{fig:para_eff} summarizes the approximation errors, averaged over 10 independent seeds, where the y-axis indicates the error and the x-axis indicates the number of trainable parameters utilized. Each curve corresponds to a fixed block size within BOFT, with the leftmost marker showing the result of BOFT($1$,$b$) (\emph{i.e.}, the standard block-diagonal OFT with block size $b$). BOFT consistently achieves superior parameter efficiency compared to OFT. For instance, BOFT($9$,$2$) outperforms BOFT($1$,$16$) in terms of expressivity, while utilizing significantly fewer parameters. In general, smaller block sizes $b$ paired with larger values of $m$ result in greater parameter efficiency; for example, BOFT($6$,$4$) requires fewer parameters while achieving comparable error to BOFT($2$,$16$). Overall, the butterfly structure imposes a richer, more organized structure within the orthogonal group than the block-diagonal approach, thereby making BOFT provably more expressive than OFT with an equivalent block size.

\begin{theorem}[Expressivity of BOFT]\label{thm:exp_boft}
    BOFT is strictly more expressive than OFT for the same block size. To approximate any orthogonal matrix of size $d$ using the butterfly structure, one can employ a product of butterfly matrices of the form $\bm{B}_{d-1,1}(d)\bm{B}^\top_{d-1,2}(d)\cdots\bm{B}_{1,1}(d)\bm{B}^\top_{1,2}(d)$, where each $\bm{B}_{i,j}(d)$ represents a butterfly matrix.
\end{theorem}

This result in Theorem~\ref{thm:exp_boft} inspires a straightforward extension for BOFT: the general orthogonal transformation can be composed as $\thickmuskip=2mu \medmuskip=2mu \bm{R}^G(m_1,b_1,m_2,b_2,l)=\bm{R}_{l,1}(m_1,b_1)\bm{R}_{l,2}^T(m_2,b_2)\cdots\bm{R}_{1,1}(m_1,b_1)\bm{R}_{1,2}^T(m_2,b_2)$, where each $\bm{R}_{i,j}^T(m,b)$ is an orthogonal matrix component as constructed in BOFT. Setting $m_1=m_2=\log d$, $b_1=b_2=2$, and $l=d-1$ enables $\bm{R}^G(m_1,b_1,m_2,b_2,l)$ to span the entire orthogonal group. This construction is also recognized as the kaleidoscope hierarchy~\cite{dao2019kaleidoscope}. Notably, while increased expressiveness enhances representational power, it does not always correlate with improved finetuning performance—unrestricted finetuning often produces suboptimal results despite its capacity to approximate any transformation. Therefore, achieving a balance between expressiveness and regularization is crucial for the success of model finetuning. By imposing structural priors, BOFT increases the parameter space accessible during finetuning and provides a more effective compromise between these two competing desiderata.

\textbf{Spectral properties}. Orthogonal finetuning is typically superior to LoRA regarding spectral characteristics, as it exactly maintains the spectral norm of the pretrained weight matrix $\bm{W}^0$. This can be illustrated using singular value decomposition: $\bm{W}^0=\bm{U}\bm{\Sigma}\bm{V}^\top$, where $\bm{U}$ and $\bm{V}$ are orthogonal matrices, and $\bm{\Sigma}$ is diagonal containing the singular values. When either OFT or BOFT is applied, they post-multiply by an orthogonal matrix $\bm{R}$ on the left, leading to the finetuned weights $\bm{R}\bm{U}\bm{\Sigma}\bm{V}^\top$. This operation preserves the largest singular value (\emph{i.e.}, the spectral norm of $\bm{W}^0$). Preservation of the spectral norm has been empirically demonstrated to be beneficial for both training stability and model generalization~\cite{miyato2018spectral, yoshida2017spectral}. Additional mathematical insights can be found in Appendix~\ref{app:sn_property}.

\textbf{Orthogonal finetuning as learning bilinear similarity}. BOFT is interpretable as a process of learning bilinear similarities of the form $\bm{w}^0_i\bm{R}\bm{x}$, where $\bm{w}^0_i$ denotes the $i$th neuron (i.e., a column of $\bm{W}^0$). In this framework, BOFT focuses on discovering an orthogonal similarity matrix $\bm{R}$, imposing a strong structural regularization, which closely ties to distance metric learning~\cite{xing2002distance} and the mathematical theory of bilinear forms~\cite{roman2005advanced}.

\textbf{Inductive bias and generalization in BOFT}. In BOFT, $\bm{R}(m,b)$ generally belongs to a structured family of the orthogonal group, restricting the hypothesis space, and thus introduces a built-in inductive bias. The inductive bias conferred by the butterfly factorization is hypothesized to be advantageous for generalization, as it encodes shared structural regularities present in classical linear operations~\cite{dao2019kaleidoscope}, such as the discrete Fourier, discrete sine/cosine, and Hadamard transforms. Furthermore, the use of sparse matrix factorization in BOFT is likely to introduce additional implicit inductive biases~\cite{gunasekar2017implicit,li2018algorithmic,liu2019neural}.

\textbf{Comparison to butterfly-based sparse training}. Numerous studies~\cite{dao2019learning,dao2019kaleidoscope,chen2022pixelated,dao2022monarch} have explored sparse training approaches that utilize butterfly parameterization. These methods generally involve directly reparameterizing neural network weight matrices using the butterfly parameterization and then training the models from the ground up. Notably, \cite{dao2022monarch} investigates finetuning of pretrained weights by initially projecting them onto a modified set of butterfly matrices, followed by optimization of the resulting components tailored to specific tasks. In contrast, BOFT introduces a fundamentally different finetuning methodology, where transformation of weights occurs via weight matrices that are shared across layers.

\input{rebuttal-results}

\section{Concluding Remarks and Limitations}

This work introduces Orthogonal Butterfly, a versatile and parameter-efficient finetuning approach tailored for foundation models and constructed on the butterfly framework. Central to achieving improved parameter-efficiency is the representation of a dense orthogonal matrix as a product of several sparse orthogonal matrices. To facilitate the identification of practical matrix factorizations, we present a graphical information transmission paradigm. Within this construct, we demonstrate that the butterfly structure meets the criteria necessary for effective sparse orthogonal matrix decomposition. Through extensive experiments, we validate that BOFT is effective for finetuning across diverse model types, including large language models, large vision models, and text-to-image generative architectures. Our empirical findings highlight BOFT's capability as a general-purpose method for model finetuning.

Nevertheless, BOFT does have certain shortcomings. Because it forms the final orthogonal matrix as a sequence of multiplications among several orthogonal matrices, the training process incurs somewhat greater runtime overhead compared to OFT. Addressing the training efficiency of BOFT is an area that remains unresolved. Importantly, post-finetuning, the orthogonal matrices derived through BOFT can be incorporated directly into the pretrained models, thus avoiding any additional latency at inference. Additionally, the optimality of the butterfly network for information transmission remains unproven. Our proposed information transmission framework also opens the door to insights from computer networking research, a field that rigorously examines network topology efficiency for information dissemination. This raises the prospect that even more efficient network architectures could potentially be leveraged for the composition of dense orthogonal matrices.